\documentclass[11pt]{article}
% ---- style ----
\usepackage[utf8]{inputenc}
\usepackage[T1]{fontenc}
\usepackage{lmodern}
\usepackage{microtype}
\usepackage{amssymb}
\usepackage[a4paper,margin=0.85in]{geometry}
\usepackage{xcolor}
\definecolor{Accent}{HTML}{1565C0}
\usepackage{titlesec}
\titleformat{\section}{\normalsize\bfseries\color{Accent}}{}{0em}{}
\titlespacing*{\section}{0pt}{6pt}{2pt}
\usepackage{enumitem}
\setlist{nosep,leftmargin=1.2em}
\usepackage[colorlinks=true,urlcolor=Accent,linkcolor=Accent]{hyperref}
\input{glyphtounicode}\pdfgentounicode=1
\setlength{\parindent}{0pt}
\setlength{\parskip}{4pt}
\pagestyle{plain}
% ---------------
\begin{document}

\begin{center}
{\large\bfseries\color{Accent} Research Statement: Past and Current Research}\\[2pt]
Alexander Salganik, University of Turku
\end{center}

Accreting X-ray pulsars let us study plasma physics in magnetic fields of
$10^{12}$--$10^{13}$\,G, across five orders of magnitude in luminosity. My research
measures the geometry, magnetic field and dynamics of their accretion flows
observationally, combining broadband spectroscopy, timing and X-ray polarimetry. I lead
my own observing programmes: I am lead author of ten papers and PI of 11 approved
NuSTAR, IXPE, Swift and Insight-HXMT programmes, including $\sim$630\,ks of Director's
Discretionary Time, and of many NICER and Swift target-of-opportunity campaigns.

\section*{Magnetic field strength and orientation: cyclotron lines and polarimetry}
Cyclotron resonant scattering features (CRSFs) measure the strength of the magnetic
field; X-ray polarization measures its orientation. I work with both. As PI of joint
IXPE and NuSTAR observations, I obtained the \textbf{first detection of X-ray
polarization from a long-period X-ray pulsar}, 4U~1954$+$319 (A\&A 710, A36), showing
that IXPE can constrain the geometry even of slowly rotating, low-luminosity systems. I
now lead IXPE programmes on GX~1$+$4 and 1E~1145.1$-$6141, contribute to IXPE studies of
X-ray pulsars, neutron-star low-mass X-ray binaries and the black-hole binary Cyg~X-1,
and, in an invited ISSI talk, tested methods to recover pulsar geometry from simulated
polarization data. On the cyclotron side, I discovered a \textbf{pulse-phase-transient
CRSF} in Swift~J1808.4$-$1754 (MNRAS 514, 2707), tracked a stable $\sim$55\,keV CRSF
through the 2026 outburst of 1A~1118$-$61, and lead an Insight-HXMT programme on
cyclotron-line harmonics.

\section*{The accretion flow has a memory}
Using dense SRG/ART-XC, NuSTAR and Insight-HXMT monitoring of 1A~1118$-$61, I discovered
a \textbf{new phenomenon: frequency hysteresis of quasi-periodic oscillations} in an
X-ray pulsar (arXiv:2611.20091). The QPO frequencies are consistent with Keplerian
motion near the inner disc edge, yet at the same luminosity they are $\sim$35\% higher on
the rise of the outburst than on the decline. The disc--magnetosphere coupling therefore
depends on accretion history, an effect not included in standard models of
magnetospheric accretion. Earlier, I discovered a 0.8-mHz QPO and abrupt timing
transitions in SXP\,31.0 (A\&A 705, A141), based on NuSTAR and NICER campaigns I led.

\section*{Accretion at low and extreme rates}
In the persistent Be/X-ray pulsar 2RXP~J130159.6$-$635806, after two decades of spin-up
the source reached spin equilibrium and switched, within days, into a \textbf{new stable
low-luminosity state} ten times fainter (A\&A 698, A71). Persistent pulsations and a
two-hump spectrum show that accretion continues, ruling out a propeller and pointing to
a field of $\gtrsim$$10^{13}$\,G. In 1A~1118$-$61 I traced how the spectrum develops a
distinct high-energy hump, containing the CRSF, as the source fades, constraining the
transition to $L\simeq(0.8$--$1.8)\times10^{36}$\,erg\,s$^{-1}$ (A\&A, in press). At the
opposite extreme, I showed that RX~J0440.9$+$4431 entered the supercritical regime
(MNRAS 524, 5213; IKI RAS Best Scientific Work Prize).

\section*{Looking ahead}
These results point to one goal: a self-consistent picture of magnetized accretion in
which field strength (from CRSFs), field orientation and emission geometry (from
polarimetry) and accretion history (from timing) are measured together across regimes.
Extending this to the soft X-ray band, which MKI's REDSoX and GOSoX missions will open,
is the focus of my research proposal.

\end{document}
